\section{Exact \texorpdfstring{\(2r\)}{2r} clusters for separated interfaces}
\label{sec:exact-2r}

\subsection{Operators and the one-interface input}

Let \(Q^{(6)}\in\{\pm1\}^{\Z}\) be the flux word whose positive sites are
\[
 \{ -4j:j\geq 0\}\cup\{6+4j:j\geq0\}.
\]
Choose a lift \(\tau^{(6)}\) with
\(\tau^{(6)}_{i+1}=Q^{(6)}_i\tau^{(6)}_i\), and define
\begin{equation}
 (A_6u)_i=u_{i-1}+u_{i+1}
 +\tau^{(6)}_{i-2}u_{i-2}+\tau^{(6)}_iu_{i+2},
 \qquad H_6=A_6^2 .
 \label{eq:g6-operator}
\end{equation}
The other lift is \(-\tau^{(6)}\); after squaring it is diagonally unitarily
equivalent to the displayed operator.  Reflection gives the reverse
orientation.

Define
\begin{align}
 p_6(y)={}&16y^{10}-520y^9+6913y^8-48448y^7+191768y^6
 -423904y^5 \notag\\
 &+484528y^4-270464y^3+137856y^2-19968y+256.
 \label{eq:p6}
\end{align}
The constant \(\cSix\) is the unique root of \(p_6\) in the rational
interval
\begin{equation}
 \frac{7905369311620327}{10^{15}}<\cSix<
 \frac{7905369311620328}{10^{15}}.
 \label{eq:c6-interval}
\end{equation}
The period-eight reference squared edge is
\begin{equation}
 \etaRef=4+\sqrt{10+2\sqrt5}<\cSix.
 \label{eq:eta}
\end{equation}

The proof below uses the following certified one-interface and Floquet
inputs.  They are the only computer-assisted spectral inputs to this
section:
\begin{align}
 \ker(H_6-\cSix)&=\Span\{\psi_+,\psi_-\},
 &\dim\ker(H_6-\cSix)&=2, \label{eq:rank-two-input}\\
 H_6\big|_{\ker(H_6-\cSix)^\perp}&\leq \cSix-\frac1{100},
 \label{eq:complement-input}
\end{align}
where the modes are orthonormal and satisfy
\(A_6\psi_\pm=\pm\sqrt{\cSix}\,\psi_\pm\).  For all eight cuts of the
period-eight bulk, both stable multipliers have modulus below
\begin{equation}
 \qF=\frac9{25},
 \label{eq:qf}
\end{equation}
and the selected two-column Floquet bases have condition number below
\(17\).  These bounds hold for both tails and are uniform under reflection.

\subsection{Finite-ring geometry and patch identification}

Work on \(\Z/n\Z\) in Hamilton gauge.  All step-one signs are positive
except for a movable seam carrying a holonomy
\(\alpha\in\{+1,-1\}\).  The flux word is
\(Q_i=\tau_i\tau_{i+1}\).  Assume that its positive sites have cyclic gaps
four except at exactly \(r\in\{1,2,3\}\) marked gaps of length six.  Write
\(p_j\) for the left positive endpoint of the \(j\)-th marked gap, ordered
cyclically, and put
\begin{equation}
 d_j=p_{j+1}-p_j\pmod n,\qquad D=\min_j d_j,
 \label{eq:separation}
\end{equation}
with \(p_{r+1}=p_1\).  For \(r=1\), the return arc has length \(d_1=D=n\).
Thus \(D\) is a distance in sites, not a number of period-eight cells.
Assume from now on that
\begin{equation}
 D\geq1040,\qquad S=\left\lfloor\frac D4\right\rfloor,
 \qquad a=S-8.
 \label{eq:S}
\end{equation}

On the oriented arc from \(p_j\) to \(p_{j+1}\), use the coordinate
\(0\leq t\leq d_j\).  Define interface cutoffs \(\chi_j\) and a bulk cutoff
\(\beta_j\) by
\begin{equation}
\begin{array}{c|cc}
\text{range of }t&\chi_j&\beta_j\\ \hline
0\leq t\leq a&1&0\\
a<t<a+S&\cos\vartheta&\sin\vartheta\\
a+S\leq t\leq d_j-a-S&0&1
\end{array}
\qquad
 \vartheta=\frac{\pi(t-a)}{2S},
 \label{eq:left-cutoff}
\end{equation}
and on the right transition by
\begin{equation}
 \beta_j=\cos\varphi,\qquad \chi_{j+1}=\sin\varphi,\qquad
 \varphi=\frac{\pi(t-d_j+a+S)}{2S}
 \label{eq:right-cutoff}
\end{equation}
for \(d_j-a-S<t<d_j-a\), followed by \(\chi_{j+1}=1\) on
\(d_j-a\leq t\leq d_j\).  Endpoint values agree.  For \(r=1\), the two
endpoint pieces are the same cyclic interface cutoff.  Pointwise,
\begin{equation}
 \sum_{j=1}^r\chi_j^2+\sum_{j=1}^r\beta_j^2=1.
 \label{eq:partition}
\end{equation}
The pure-bulk plateau on the \(j\)-th arc is
\([p_j+2S-8,p_{j+1}-2S+8]\), and its length is at least
\begin{equation}
 d_j-4S+16\geq D-4\lfloor D/4\rfloor+16\geq16.
 \label{eq:bulk-plateau}
\end{equation}

For a set \(E\), let \(E^{[k]}\) denote its cyclic \(k\)-neighborhood.
Set
\[
 I_j=\{i:\chi_j(i)\neq0\},\quad J_j=I_j^{[4]},\qquad
 B_j=\{i:\beta_j(i)\neq0\},\quad C_j=B_j^{[4]}.
\]
Since \(H=A^2\) has range four, the coefficient collar needed to identify a
compression to \(J_j\) is contained in \(I_j^{[6]}\).  That collar stays at
distance at least
\begin{equation}
 D-2\lfloor D/4\rfloor-4\geq516
 \label{eq:collar-separation}
\end{equation}
from every other marked core.  Likewise \(B_j^{[6]}\) stays at distance at
least \(S-19\geq241\) from each six-gap core.  Hence every interface patch
contains exactly one marked core and every bulk patch is a pure
period-eight patch.

The holonomy seam can be moved by diagonal switching.  Place it on an edge
of a bulk plateau whose endpoints have distance at least eight from every
interface support.  It then misses every \(I_j^{[6]}\).  Each interface
collar is a proper arc, so after unwrapping it, translation, optional
reflection, and diagonal switching identify its coefficients exactly with
the corresponding finite section of \(H_6\).  On a proper arc, two lifts of
the same \(Q\)-word differ by a global sign.  If
\((D_0u)_i=(-1)^iu_i\), then
\begin{equation}
 A_{-\tau}=-D_0A_\tau D_0,\qquad H_{-\tau}=D_0H_\tau D_0.
 \label{eq:lift-equivalence}
\end{equation}
This removes the remaining lift choice after squaring.  The same argument
identifies each \(C_j\) with a finite section of the period-eight reference
operator.  We record the result in the form needed below.

\begin{lemma}[Patch identification]
\label{lem:patch-identification}
For every allowed orientation, lift, and holonomy, there are local unitaries
such that:
\begin{enumerate}[label=\textup{(\roman*)}]
\item the compression of \(H\) to \(J_j\) is a principal finite section of
the corresponding infinite G6 operator;
\item the image under \(H\) of a vector supported in \(I_j\) is supported in
\(J_j\) and agrees with the image under that infinite operator;
\item the analogous compression and map on \(B_j\subset C_j\) agree with a
period-eight reference-bulk finite section.
\end{enumerate}
These statements include cyclic wraparound because every collar is first
unwrapped as a proper arc.
\end{lemma}

Transport the orthonormal pair \(\psi_+,\psi_-\) through the same local
unitary, before truncation, and call the transported modes
\(\psi_{j,+},\psi_{j,-}\).  For every finite-ring vector \(x\), zero
extension on the unwrapped arc gives
\begin{equation}
 \langle\psi_{j,\pm},\chi_jx\rangle_{\ell^2(\Z)}
 =\langle\chi_j\psi_{j,\pm},x\rangle_{\C^n}.
 \label{eq:complement-identity}
\end{equation}
This identity is the bridge that permits the complete rank-two complement
bound \eqref{eq:complement-input} to be used on each localized vector.

\subsection{Tails, Gram matrix, and the lower count}

The cutoff starts changing at distance \(S-8\) from a marked core.  Removing
the propagation range four leaves
\begin{equation}
 L_{\mathrm{site}}=S-12,\qquad
 \ell=\left\lfloor\frac{L_{\mathrm{site}}}{8}\right\rfloor.
 \label{eq:ell}
\end{equation}
The eight site residues, the two tails, the condition-number bound, and the
geometric series imply, for either normalized local mode,
\begin{equation}
 \|\psi_{j,\pm}\|_{\mathrm{tail}}^2
 \leq\frac{16\cdot17^2}{1-\qF^2}\qF^{2\ell}
 =\frac{10625}{2}\qF^{2\ell}<73^2\qF^{2\ell}.
 \label{eq:tail}
\end{equation}
Put
\begin{equation}
 \phi_{j,\pm}=\chi_j\psi_{j,\pm},
 \qquad
 \Phi=(\phi_{1,+},\phi_{1,-},\ldots,\phi_{r,+},\phi_{r,-}),
 \qquad G=\Phi^*\Phi.
 \label{eq:phi}
\end{equation}
The full modes at one interface are orthogonal.  Their truncated cross term
is not discarded; rather,
\[
 |\langle\phi_{j,+},\phi_{j,-}\rangle|
 \leq73^2\qF^{2\ell}.
\]
The same estimate holds between different interfaces because both vectors
are then in their tails.  With \(m=2r\),
\begin{equation}
 \|G-\Id_m\|\leq m73^2\qF^{2\ell}.
 \label{eq:gram}
\end{equation}
The hypotheses give \(S\geq260\), \(L_{\mathrm{site}}\geq248\), and
\(\ell\geq31\).  Since \(m\leq6\),
\begin{equation}
 6\cdot73^2\left(\frac9{25}\right)^{62}<\frac12.
 \label{eq:gram-numeric}
\end{equation}
Thus the \(2r\) columns are independent and \(\|G^{-1}\|\leq2\).

The row-sum bound gives \(\|H\|\leq16\), while \(\cSix<8\).  Cutting an
exact mode only in its tails and using \eqref{eq:tail} yields
\begin{equation}
 \|(H-\cSix)\phi_{j,\pm}\|
 \leq(16+8)73\qF^\ell=1752\qF^\ell.
 \label{eq:residual-column}
\end{equation}
Define the isometry
\begin{equation}
 U=\Phi G^{-1/2}:\C^{2r}\longrightarrow\C^n.
 \label{eq:U}
\end{equation}
Then
\begin{equation}
 \|(H-\cSix)U\|<3504r\qF^\ell<\frac1{400}.
 \label{eq:residual-U}
\end{equation}
Let \(E_W\) be the spectral projection of \(H\) onto
\begin{equation}
 W=\left[\cSix-\frac1{400},\cSix+\frac1{400}\right].
 \label{eq:window}
\end{equation}
If \(\rank E_W<2r\), a nonzero vector in \(\ran U\) is orthogonal to
\(\ran E_W\).  The spectral theorem then makes its residual at least
\(1/400\), contradicting \eqref{eq:residual-U}.  Therefore
\begin{equation}
 \rank E_W\geq2r.
 \label{eq:lower-count}
\end{equation}

\subsection{The codimension-\texorpdfstring{\(2r\)}{2r} complement}

Let \(V=\ran\Phi\).  If \(x\perp V\), then
\eqref{eq:complement-identity} gives
\[
 \langle\psi_{j,+},\chi_jx\rangle
 =\langle\psi_{j,-},\chi_jx\rangle=0
\]
for every interface.  Thus each interface-localized vector is orthogonal to
the entire two-dimensional space in \eqref{eq:rank-two-input}, and
\eqref{eq:complement-input} applies.  Each bulk-localized vector satisfies
the stronger upper bound \(\etaRef<\cSix-1/100\).

The cutoff vector follows a unit-circle quarter arc at angular speed
\(\pi/(2S)\).  If two sites have cyclic distance \(h\leq4\), the chord bound
gives
\[
 \sum_\gamma|\gamma(u)-\gamma(v)|^2
 \leq\frac{\pi^2h^2}{4S^2},
\]
where \(\gamma\) ranges over all cutoffs.  The exact double-commutator IMS
identity, the range-four row-sum bound, and \(\pi^2<10\) consequently give
\begin{equation}
 \|E_{\mathrm{IMS}}\|<\frac{320}{S^2}
 \leq\frac{320}{260^2}=\frac4{845}.
 \label{eq:ims}
\end{equation}
Using \eqref{eq:partition} and summing the localized quadratic-form bounds,
we obtain on \(V^\perp\)
\begin{equation}
 H\big|_{V^\perp}
 \leq\cSix-\frac1{100}+\frac4{845}
 =\cSix-\frac{89}{16900}
 <\cSix-\frac1{200}.
 \label{eq:complement-finite}
\end{equation}
By the codimension form of min--max, at most \(\dim V=2r\) eigenvalues lie
above \(\cSix-1/200\).  Since the entire window \(W\) lies above this
level, \eqref{eq:lower-count} sharpens to
\begin{equation}
 \rank E_W=2r.
 \label{eq:exact-count}
\end{equation}
The equality counts algebraic multiplicity and makes no assertion of
individual simplicity.

\subsection{The \texorpdfstring{\(2r\)}{2r}-dimensional Feshbach equation}

Let
\begin{equation}
 P=UU^*,\qquad \perpQ=\Id-P,\qquad E=(H-\cSix)\Phi.
 \label{eq:PE}
\end{equation}
Because \(\ran U=V\), estimate \eqref{eq:complement-finite} holds on
\(\ran\perpQ\).  If \(|z-\cSix|\leq1/400\), then
\begin{equation}
 \|\big(\perpQ H\perpQ-z\big)^{-1}\|\leq400
 \quad\text{on }\ran\perpQ.
 \label{eq:resolvent}
\end{equation}
Block Gaussian elimination with respect to \(P+\perpQ=\Id\) gives the
effective operator on \(\C^{2r}\),
\begin{equation}
 H_{\mathrm{eff}}(z)=U^*HU-U^*H\perpQ
 (\perpQ H\perpQ-z)^{-1}\perpQ HU,
 \label{eq:feshbach}
\end{equation}
and the finite-ring eigenvalue equation in \(W\) is exactly
\begin{equation}
 \det\big(H_{\mathrm{eff}}(z)-z\Id_{2r}\big)=0.
 \label{eq:feshbach-determinant}
\end{equation}
Using \(U=\Phi G^{-1/2}\) and
\(\perpQ HU=\perpQ EG^{-1/2}\), one obtains
\begin{align}
 H_{\mathrm{eff}}(z)-\cSix\Id_{2r}
={}&G^{-1/2}\Phi^*EG^{-1/2} \notag\\
 &-G^{-1/2}E^*\perpQ(\perpQ H\perpQ-z)^{-1}
 \perpQ EG^{-1/2}.
 \label{eq:feshbach-expanded}
\end{align}
No orthogonality of the uncorrected columns of \(\Phi\) is assumed.
Writing the right side as \(T_1+R_2(z)\), the Gram and residual bounds give
\begin{equation}
 \|T_1\|\leq3504r\qF^\ell,\qquad
 \|R_2(z)\|\leq400r\,3504^2\qF^{2\ell}.
 \label{eq:feshbach-bounds}
\end{equation}
Since
\begin{equation}
 400\cdot3504^2\left(\frac9{25}\right)^{31}<1,
 \label{eq:remainder-numeric}
\end{equation}
we have \(\|R_2(z)\|<r\qF^\ell\).  Every root of
\eqref{eq:feshbach-determinant} in \(W\) therefore satisfies
\begin{equation}
 |\lambda_j-\cSix|<3505r\left(\frac9{25}\right)^\ell,
 \qquad 1\leq j\leq2r.
 \label{eq:cluster-bound}
\end{equation}

\begin{theorem}[Separated-interface cluster]
\label{thm:exact-2r}
Let a legal finite ring contain exactly \(r\in\{1,2,3\}\) elementary G6
interfaces, with arbitrary interface orientations, either lift, and either
Hamilton holonomy.  If their minimum cyclic site separation satisfies
\(D\geq1040\), define \(S,L_{\mathrm{site}},\ell\) by
\eqref{eq:S} and \eqref{eq:ell}.  Then
\[
 \rank\one_{[\cSix-1/400,\,\cSix+1/400]}(H)=2r.
\]
The \(2r\) eigenvalues in this window, counted with multiplicity, satisfy
\eqref{eq:cluster-bound}, and their exact finite-ring equation is the
\(2r\)-dimensional determinant \eqref{eq:feshbach-determinant}.
\end{theorem}

\subsection{A sufficient exponential onset}

The explicit constants also give a continuous sufficient onset for the
exponential construction.  For every even \(n\),
\[
\rho_-(n)^2>8-\frac{200}{n^2}.
\]
For the nonzero residues put
\[
B(n)=8-\frac{200}{n^2}-\frac9{100}.
\]
If the corresponding residue construction has \(r\) interfaces and minimum
separation \(D\geq1040\), the cluster bound gives
\[
\rho(A_n)^2
\leq c_6^{\mathrm{up}}
+3505r\left(\frac9{25}\right)^\ell,
\qquad
\ell=\left\lfloor
\frac{\lfloor D/4\rfloor-12}{8}
\right\rfloor,
\]
where \(c_6^{\mathrm{up}}=7905369311620328/10^{15}\).
The first distance-eligible endpoints are
\[
\begin{array}{c|c|c|c|c}
n\bmod8&n&r&D&\ell\\ \hline
2&1042&1&1042&31\\
4&2084&2&1042&31\\
6&3126&3&1042&31
\end{array}
\]
and exact rational evaluation gives
\[
B(n)-\left[
c_6^{\mathrm{up}}
+3505r\left(\frac9{25}\right)^{31}
\right]>\frac1{250}
\]
in all three rows.  Along each residue subsequence, \(D\) and \(\ell\) do
not decrease, the spectral cap does not increase, and \(B(n)\) strictly
increases.

For residue zero, the period-eight cap is \(1561/200\).  At \(n=3120\),
\[
8-\frac{200}{3120^2}=\frac{389375}{48672},
\qquad
\frac{389375}{48672}-\frac{1561}{200}
=\frac{237251}{1216800}>0.
\]
The first four even orders at or above \(3120\) are \(3120,3122,3124,3126\),
one in each residue class.  The preceding endpoint checks and residue-wise
monotonicity therefore prove a strict counterexample for every even
\(n\geq3120\).  We record
\[
N_{\mathrm{exp}}=3120
\]
as a sufficient, nonoptimal onset for this exponential estimate.  It is not
the sharp onset: the finite bridge plus IMS theorem in the main paper gives
the stronger value \(48\).

The restrictions \(r\leq3\) and \(D\geq1040\) enter, respectively, in the
uniform Gram estimate and in the seam-clearance, tail, and IMS bounds.  The
argument proves neither individual simplicity nor a nonzero leading
interaction coefficient.
